% GENERATED FILE. Edit canonical lessons, then run npm run generate:books.
% canonical-source-hash: fnv1a64:119051b44cd49762
% canonical-lessons: ML-C427-daridryam, ML-C427-daridran, ML-C427-pattini, ML-C427-ceri, ML-C427-saksarata

\chapter{Poverty}
\label{ch:ml-poverty}
\hlchaptermodality{427}

\begin{chapteropening}
\textbf{By the end of this chapter:} \emph{I can say poverty, a poor man, hunger, a slum and literacy.}

Complete the last lesson of chapter 427: I can say poverty, a poor man, hunger, a slum and literacy.
\end{chapteropening}

\section[\texorpdfstring{\ml{ദാരിദ്ര്യം} (dāridryaṁ)}{ദാരിദ്ര്യം (dāridryaṁ)}]{\ml{ദാരിദ്ര്യം} (dāridryaṁ) — poverty}

\label{lesson:ML-C427-daridryam}



\begin{quote}
Before the new one: say the Malayalam for to submit, then the Malayalam for a petition.
\end{quote}

\subsection*{You'll want to know: \ml{ദാരിദ്ര്യം}}
\textbf{\ml{ദാരിദ്ര്യം}} — \emph{dāridryaṁ} — \textquotedblleft{}poverty\textquotedblright{}.

\begin{quote}
\ml{ദാരിദ്ര്യം} \ml{കാരണം} \ml{ആ} \ml{കുടുംബം} \ml{നഗരത്തിലേക്ക്} \ml{മാറി}.
\end{quote}

\begin{quote}
\emph{dāridryaṁ kāraṇaṁ ā kuṭumbaṁ nagarattilēkkŭ māṟi.}
\end{quote}

\begin{quote}
Because of poverty, that family moved to the city.
\end{quote}

Below the poverty line: \ml{ദാരിദ്ര്യരേഖയ്ക്ക്} \ml{താഴെ}.

\begin{cousinweb}
\textbf{\ml{ദാരിദ്ര്യം}} (\emph{dāridryaṁ}) is borrowed from Sanskrit \emph{dāridrya}, \textquotedblleft{}poverty\textquotedblright{} (CDIAL 6297), made from \emph{daridra}, \textquotedblleft{}poor\textquotedblright{}.
\end{cousinweb}

\begin{grammarlens}[title={What you've built}]
One more word for giving reasons and opinions.
\end{grammarlens}

\subsection*{Your turn}
\begin{itemize}
\raggedright
  \item \emph{Say it:} \emph{dāridryaṁ}
  \item \emph{Say it:} \emph{dāridryaṁ}, once more
  \item \emph{Say it:} \emph{dāridryaṁ kāraṇaṁ ā kuṭumbaṁ nagarattilēkkŭ māṟi.} — Because of poverty, that family moved to the city.
  \item \emph{Recall:} say the Malayalam for to submit, then the Malayalam for a petition, then say \emph{dāridryaṁ} again
\end{itemize}

\begin{tcolorbox}[breakable,colback=teal!4,colframe=teal!35!black,title={Before you move on}]
What does \ml{ദാരിദ്ര്യം} mean? (\textquotedblleft{}Poverty\textquotedblright{}.) Say it once more.
\end{tcolorbox}
\section[\texorpdfstring{\ml{ദരിദ്രൻ} (daridran)}{ദരിദ്രൻ (daridran)}]{\ml{ദരിദ്രൻ} (daridran) — a poor man}

\label{lesson:ML-C427-daridran}



\begin{quote}
Before the new one: say the Malayalam for a petition, then the Malayalam for poverty.
\end{quote}

\subsection*{You'll want to know: \ml{ദരിദ്രൻ}}
\textbf{\ml{ദരിദ്രൻ}} — \emph{daridran} — \textquotedblleft{}a poor man\textquotedblright{}.

\begin{quote}
\ml{ആ} \ml{ദരിദ്രന്} \ml{ഗ്രാമം} \ml{ഒരു} \ml{വീട്} \ml{കൊടുത്തു}.
\end{quote}

\begin{quote}
\emph{ā daridranŭ grāmaṁ oru vīṭŭ koṭuttu.}
\end{quote}

\begin{quote}
The village gave that poor man a house.
\end{quote}

Plural: \ml{ദരിദ്രർ}, \emph{daridrar}.

Plural: \ml{ദരിദ്രർ}, the poor. A poor woman is \ml{ദരിദ്ര}.

\begin{cousinweb}
\textbf{\ml{ദരിദ്രൻ}} (\emph{daridran}) is borrowed from Sanskrit \emph{daridra}, \textquotedblleft{}roving\textquotedblright{}, and later \textquotedblleft{}poor\textquotedblright{} (CDIAL 6195), with the Malayalam masculine ending -\ml{ൻ}.
\end{cousinweb}

\begin{grammarlens}[title={What you've built}]
One more word for giving reasons and opinions.
\end{grammarlens}

\subsection*{Your turn}
\begin{itemize}
\raggedright
  \item \emph{Say it:} \emph{daridran}
  \item \emph{Say it:} \emph{daridran}, once more
  \item \emph{Say it:} \emph{ā daridranŭ grāmaṁ oru vīṭŭ koṭuttu.} — The village gave that poor man a house.
  \item \emph{Recall:} say the Malayalam for a petition, then the Malayalam for poverty, then say \emph{daridran} again
\end{itemize}

\begin{tcolorbox}[breakable,colback=teal!4,colframe=teal!35!black,title={Before you move on}]
What does \ml{ദരിദ്രൻ} mean? (\textquotedblleft{}A poor man\textquotedblright{}.) Say it once more.
\end{tcolorbox}
\section[\texorpdfstring{\ml{പട്ടിണി} (paṭṭiṇi)}{പട്ടിണി (paṭṭiṇi)}]{\ml{പട്ടിണി} (paṭṭiṇi) — starvation; hunger}

\label{lesson:ML-C427-pattini}



\begin{quote}
Before the new one: say the Malayalam for poverty, then the Malayalam for a poor man.
\end{quote}

\subsection*{You'll want to know: \ml{പട്ടിണി}}
\textbf{\ml{പട്ടിണി}} — \emph{paṭṭiṇi} — \textquotedblleft{}starvation; hunger\textquotedblright{}.

\begin{quote}
\ml{ജോലി} \ml{ഇല്ലാതെ} \ml{ആ} \ml{കുടുംബം} \ml{പട്ടിണിയിലായി}.
\end{quote}

\begin{quote}
\emph{jōli illāte ā kuṭumbaṁ paṭṭiṇiyilāyi.}
\end{quote}

\begin{quote}
Without work, that family fell into hunger.
\end{quote}

To go hungry is \ml{പട്ടിണി} \ml{കിടക്കുക}, to lie in hunger.

\begin{cousinweb}
\textbf{\ml{പട്ടിണി}} (\emph{paṭṭiṇi}) is native Dravidian: DEDR 3872 lists Malayalam \emph{paṭṭiṇi}, \textquotedblleft{}privation of food, starvation\textquotedblright{}, with Tamil \emph{paṭṭiṉi}, \textquotedblleft{}fasting\textquotedblright{}.
\end{cousinweb}

\begin{grammarlens}[title={What you've built}]
One more word for giving reasons and opinions.
\end{grammarlens}

\subsection*{Your turn}
\begin{itemize}
\raggedright
  \item \emph{Say it:} \emph{paṭṭiṇi}
  \item \emph{Say it:} \emph{paṭṭiṇi}, once more
  \item \emph{Say it:} \emph{jōli illāte ā kuṭumbaṁ paṭṭiṇiyilāyi.} — Without work, that family fell into hunger.
  \item \emph{Recall:} say the Malayalam for poverty, then the Malayalam for a poor man, then say \emph{paṭṭiṇi} again
\end{itemize}

\begin{tcolorbox}[breakable,colback=teal!4,colframe=teal!35!black,title={Before you move on}]
What does \ml{പട്ടിണി} mean? (\textquotedblleft{}Starvation; hunger\textquotedblright{}.) Say it once more.
\end{tcolorbox}
\section[\texorpdfstring{\ml{ചേരി} (cēri)}{ചേരി (cēri)}]{\ml{ചേരി} (cēri) — a slum}

\label{lesson:ML-C427-ceri}



\begin{quote}
Before the new one: say the Malayalam for a poor man, then the Malayalam for starvation.
\end{quote}

\subsection*{You'll want to know: \ml{ചേരി}}
\textbf{\ml{ചേരി}} — \emph{cēri} — \textquotedblleft{}a slum\textquotedblright{}.

\begin{quote}
\ml{നഗരത്തിലെ} \ml{ചേരിയിൽ} \ml{ആയിരക്കണക്കിന്} \ml{ആളുകൾ} \ml{താമസിക്കുന്നു}.
\end{quote}

\begin{quote}
\emph{nagarattile cēriyil āyirakkaṇakkinŭ āḷukaḷ tāmasikkunnu.}
\end{quote}

\begin{quote}
Thousands of people live in the city's slum.
\end{quote}

Plural: \ml{ചേരികൾ}.

\begin{cousinweb}
\textbf{\ml{ചേരി}} (\emph{cēri}) is native Dravidian: DEDR 2007 lists Malayalam \emph{cēri}, \textquotedblleft{}an assemblage, a village street\textquotedblright{}, with Tamil \emph{cēri}, \textquotedblleft{}a hamlet\textquotedblright{}. The sense \textquotedblleft{}slum\textquotedblright{} is the modern one.
\end{cousinweb}

\begin{grammarlens}[title={What you've built}]
One more word for giving reasons and opinions.
\end{grammarlens}

\subsection*{Your turn}
\begin{itemize}
\raggedright
  \item \emph{Say it:} \emph{cēri}
  \item \emph{Say it:} \emph{cēri}, once more
  \item \emph{Say it:} \emph{nagarattile cēriyil āyirakkaṇakkinŭ āḷukaḷ tāmasikkunnu.} — Thousands of people live in the city's slum.
  \item \emph{Recall:} say the Malayalam for a poor man, then the Malayalam for starvation, then say \emph{cēri} again
\end{itemize}

\begin{tcolorbox}[breakable,colback=teal!4,colframe=teal!35!black,title={Before you move on}]
What does \ml{ചേരി} mean? (\textquotedblleft{}A slum\textquotedblright{}.) Say it once more.
\end{tcolorbox}
\section[\texorpdfstring{\ml{സാക്ഷരത} (sākṣarata)}{സാക്ഷരത (sākṣarata)}]{\ml{സാക്ഷരത} (sākṣarata) — literacy}

\label{lesson:ML-C427-saksarata}



\begin{quote}
Before the new one: say the Malayalam for starvation, then the Malayalam for a slum.
\end{quote}

\subsection*{You'll want to know: \ml{സാക്ഷരത}}
\textbf{\ml{സാക്ഷരത}} — \emph{sākṣarata} — \textquotedblleft{}literacy\textquotedblright{}.

\begin{quote}
\ml{ആ} \ml{ഗ്രാമത്തിലെ} \ml{സാക്ഷരത} \ml{കുറവാണ്}.
\end{quote}

\begin{quote}
\emph{ā grāmattile sākṣarata kuṟavāṇŭ.}
\end{quote}

\begin{quote}
Literacy in that village is low.
\end{quote}

A literate person is \ml{സാക്ഷരൻ}.

\begin{cousinweb}
\textbf{\ml{സാക്ഷരത}} (\emph{sākṣarata}) is built from Sanskrit \emph{sākṣara}, \textquotedblleft{}containing letters\textquotedblright{} (CDIAL 13320): \emph{sa-}, \textquotedblleft{}with\textquotedblright{}, and \emph{akṣara}, \textquotedblleft{}a letter\textquotedblright{}, with the abstract ending -\ml{ത}.
\end{cousinweb}

\begin{grammarlens}[title={What you've built}]
One more word for giving reasons and opinions.
\end{grammarlens}

\subsection*{Your turn}
\begin{itemize}
\raggedright
  \item \emph{Say it:} \emph{sākṣarata}
  \item \emph{Say it:} \emph{sākṣarata}, once more
  \item \emph{Say it:} \emph{ā grāmattile sākṣarata kuṟavāṇŭ.} — Literacy in that village is low.
  \item \emph{Recall:} say the Malayalam for starvation, then the Malayalam for a slum, then say \emph{sākṣarata} again
\end{itemize}

\begin{tcolorbox}[breakable,colback=teal!4,colframe=teal!35!black,title={Before you move on}]
What does \ml{സാക്ഷരത} mean? (\textquotedblleft{}Literacy\textquotedblright{}.) Say it once more.
\end{tcolorbox}
